\documentclass[11pt,a4paper]{article}
\usepackage[T1]{fontenc}
\usepackage{isabelle,isabellesym}
\usepackage{pdfsetup}
\usepackage{amsmath}
\usepackage{amssymb}
\usepackage{url}

\isabellestyle{it}

\begin{document}
\title{Strong Normalization for Church-Style System F}
\author{Arthur Freitas Ramos \and David Barros Hulak \and Ruy J. G. B. de Queiroz}
\maketitle

\tableofcontents

\begin{abstract}
This entry formalizes strong normalization for Church-style System F with
explicit type abstraction and type application.  Types and terms use de Bruijn
indices; the reduction relation contains both term-\(\beta\) and type-\(\beta\)
steps and is closed under all term contexts.  The proof combines
Girard-style reducibility candidates~\cite{Girard1972} for an untyped
lambda-calculus erasure with an explicit type-syntax measure that reflects the
erased normalization result back to the full calculus.  The final theorem
states that every well-typed System F term is strongly normalizing, with the
so the usual closed-term result is an immediate instance.
\end{abstract}

\section{Overview}

System F extends the simply typed lambda calculus with universal types,
polymorphic abstraction, and type instantiation.  The formalization represents
these constructs explicitly:
\[
  \mathop{\mathrm{TmTAbs}}\,M
  \qquad\text{and}\qquad
  \mathop{\mathrm{TmTApp}}\,M\,A .
\]
Typing uses de Bruijn indices for both term and type binders.  The relation
\(\longrightarrow_\beta\) includes ordinary term-\(\beta\), type-\(\beta\), and
compatible-context rules.

The semantic part interprets each type as a reducibility candidate in the
untyped lambda calculus.  Universal types quantify over all candidates, so
polymorphic instantiation is handled semantically rather than by collapsing
all type variables to one simple type.  Erasure removes only type syntax and
annotations.  Every full reduction step either produces an erased
untyped-\(\beta\) step or leaves the erasure unchanged while decreasing the
number of explicit type abstractions and applications.  A lexicographic
reflection argument
therefore yields strong normalization for the original System F term.

The proof proceeds in three stages.  First, reducibility candidates establish
strong normalization of the untyped erasure of every well-typed term.  Second,
each System F reduction is shown either to induce an untyped \(\beta\)-step
after erasure or to preserve the erasure while strictly decreasing the amount
of explicit type syntax.  Finally, a lexicographic argument reflects strong
normalization of the erasure back to the Church-style calculus.

Earlier formalizations include Altenkirch's LEGO development~\cite{Altenkirch1993},
the ATS/LF development of Donnelly and Xi~\cite{DonnellyXi2007}, and the
Isabelle/HOL development of Popescu, Gunter, and Osborn~\cite{Popescu2010}.
The latter formalization uses Curry-style terms and HOAS on top of FOAS.  The
present entry instead treats Church-style syntax directly: term and type
binders use de Bruijn indices, type-\(\beta\) is an explicit reduction rule,
and the erasure proof reflects normalization with an explicit type-syntax
measure.  This distinguishes the calculus, binding representation, and proof
architecture without claiming novelty for the classical normalization theorem.

We also considered basing the development on Berghofer's
\texttt{StrongNorm} formalization in the Isabelle/HOL
\texttt{HOL-Proofs-Lambda} library~\cite{BerghoferStrongNorm}.  That
development provides a closely related de Bruijn treatment of strong
normalization for the simply typed
lambda calculus, including lifting, substitution, subject reduction, and an
inductive characterization of terminating terms.  It has only term
abstraction/application, simple arrow types, and term-\(\beta\), however; it
does not cover polymorphic type abstraction/application, impredicative
universal types, or type-\(\beta\) reduction.  The present entry therefore
follows its de Bruijn and substitution discipline while developing a
System-F-specific reducibility-candidate semantics and an erasure/reflection
argument for explicit type reductions.  Berghofer's POPLmark AFP
entry~\cite{BerghoferPOPLmark2007} is similarly valuable for the de Bruijn
treatment of binding and substitution,
but it formalizes the type-safety infrastructure for System \(F_{<:}\), not
strong normalization.  We use these developments as methodological
references rather than claiming to extend either source by direct import.

The development also proves subject reduction independently of the
normalization argument.  It is included both as a useful metatheorem and as
an independent validation of the substitution infrastructure; the
normalization proof does not invoke preservation.  Typed weakening and
term/type substitution lemmas validate the de Bruijn operations at both
term-\(\beta\) and type-\(\beta\) redexes, and the preservation theorem provides
an independent consistency check on the operational semantics.

\input{session}

\bibliographystyle{abbrv}
\bibliography{root}

\end{document}
